\documentclass[11pt]{article}
\usepackage[a4paper,margin=27mm]{geometry}
\usepackage{amsmath,amssymb,amsthm,mathtools}
\usepackage[colorlinks=true,linkcolor=blue,urlcolor=blue]{hyperref}
\newtheorem{theorem}{Theorem}
\newtheorem{lemma}[theorem]{Lemma}
\newtheorem{remark}[theorem]{Remark}
\newcommand{\F}{\mathbb F}
\newcommand{\rank}{\operatorname{rank}}
\newcommand{\SV}{\operatorname{SV}}
\newcommand{\Mat}{\operatorname{Mat}}
\title{An Affine-Orbit Consequence of the Interval-Rank Cone Theorem}
\author{Yilin Liu \and GPT-6 Astra Pro}
\date{26 October 2026}
\begin{document}
\maketitle

\noindent\textbf{Status.}
This note proves a consequence of the proposed rank-profile cone theorem
in the companion \texttt{interval-rank-cones.tex}. Its complete proof
is supplied there and is not reproduced here. The main argument
below uses that theorem; a final remark gives an alternative using its
elementary coefficient-space counterpart. The argument does not use the
width-three joint classification. It has internal model review, not
external human peer review; no novelty or priority claim is made.

\section{Model and statement}

Let $n,w,\delta\ge1$. Consider
\[
 f(X)=uA_1(\ell_1(X)+c_1)\cdots
       A_n(\ell_n(X)+c_n)v,
\]
where $u,v$ are constant endpoints, the polynomial matrices have width
at most $w$, and every entry of every univariate $A_i$ has degree at most
$\delta$. The map
\[
 L:X\longmapsto(\ell_1(X),\ldots,\ell_n(X))
\]
is invertible and linear. The map $L$, the constants $c_i$, the variable
order, and the factors are unavailable to the black-box algorithm.
Constant factors and unused coordinates are allowed. Compatible zero
padding reduces rectangular factors to square width $w$.

This is the affine orbit of the individual-degree-$\delta$, width-$w$
ROABP class. Indeed, if only the scalar preimage polynomial is promised
to have individual degree at most $\delta$, truncating each univariate
factor above that degree preserves the scalar polynomial: each variable
occurs in only its own factor.

Define
\[
 D_w(w)=w^2-1,\qquad
 D_w(r)=2r^2\left(2\max_{r<s\le w}D_w(s)+1\right)
 \quad(1\le r<w),
\]
and write
\[
 b_r=\max_{r<s\le w}D_w(s),\qquad
 k_{\rm geo}(w)=1+\sum_{r=1}^{w-1}(2b_r+2).
\]
At $w=1$, the sum is empty and $k_{\rm geo}(1)=1$.

For distinct abscissas $a_1,\ldots,a_n$ in the working field, let
$\lambda_i(y)$ be the Lagrange interpolation polynomial satisfying
$\lambda_i(a_j)=\mathbf 1_{i=j}$. Put
\[
 \SV_{h,i}(y,z)=\sum_{j=1}^h z_j\lambda_i(y_j).
\]

\begin{theorem}[Affine-orbit hitting via the companion cone theorem]\label{thm:orbit}
Use the rank-profile cone theorem stated in Section~\ref{sec:input}
and proved in the companion \texttt{interval-rank-cones.tex}.
For $1\le r<w$, put $a_r=\binom wr^2-1$ and choose
\[
 \eta_r=
 \begin{cases}
 \min\{r\delta,a_r\},&
  \operatorname{char}\F=0\text{ or }\operatorname{char}\F>w\delta,\\
 r\delta,&\text{otherwise}.
 \end{cases}
\]
Set
\[
 K=1+\sum_{r=1}^{w-1}\eta_r(2b_r+2).
\]
Then $\SV_K$ preserves the generic rank of every contiguous interval of
every matrix product in the model above, and hits every nonzero scalar
output. In particular,
\[
 K\le \delta w\,k_{\rm geo}(w),
 \qquad
 K\le \delta w\,4^{w-1}(w!)^2\quad(w\ge2).
\]
In characteristic zero, or characteristic greater than $w\delta$,
\[
 K\le4^w k_{\rm geo}(w)
   \le4^{2w-1}(w!)^2=2^{O(w\log w)}\quad(w\ge2),
\]
independently of $\delta$.
For $w=1$, $K=1$.

If the evaluation field has at least $n^2\delta+1$ elements, a grid gives
at most
\[
 (n^2\delta+1)^{2K}
\]
oracle queries. This is polynomial in $n,\delta$ for fixed $w$ in the
stated characteristic-zero or large-characteristic regime. In every
characteristic it is polynomial in $n$ for fixed $w,\delta$, including
the multilinear case $\delta=1$. Extension-field evaluations must be
permitted when the base field is too small.
\end{theorem}

\section{The input and elementary consequences}\label{sec:input}

For a polynomial matrix $M$ over a field $E$, work over an algebraic
closure $k$ of $E$ and define
\[
 \mathcal C(M)=
 \overline{\{\alpha M(x):\alpha\in k,\ x\in k^{|X|}\}}.
\]
A restriction fixing all but selected variable blocks to zero is
\emph{cone preserving} if it leaves this scalar-saturated image closure
unchanged.

\medskip
\noindent\textbf{Assumed rank-profile cone theorem.}
For a width-$w$ product on disjoint variable blocks, suppose every
contiguous interval has the same generic rank as its constant value.
A nonzero full product of rank $s$ has a cone-preserving restriction to
at most $D_w(s)$ original blocks.
\medskip

The blocks in this input may be singleton variables, which is the case
used below. The input applies over arbitrary coefficient-field
extensions.

\begin{lemma}[Cone restrictions and coefficients]\label{lem:cone}
If a matrix product $M(Y)$ has a cone-preserving restriction retaining
at most $t$ singleton variables, then the coefficient-matrix span of
$\bigwedge^r M$ is generated by coefficients of support at most $t$.
\end{lemma}
\begin{proof}
A linear functional on the entries of $\bigwedge^r M$ is homogeneous
of degree $r$ in the entries of $M$. If it vanishes on every coefficient
of the restriction, it vanishes on the restricted image cone and hence
on the full cone. The restricted and full coefficient spaces therefore
have the same annihilator and are equal. All restricted coefficients
have support at most $t$. Equality after extending to an algebraic
closure descends to the original coefficient field.
\end{proof}

\begin{lemma}[Local Taylor coefficient span]\label{lem:jets}
Let $H(T)$ be a nonzero polynomial matrix of degree at most $d$.
Write $q=\dim\operatorname{span}\operatorname{Coeff}(H)$.
In characteristic zero, or characteristic greater than $d$, there is
a nonzero polynomial $W_H(C)$ such that, whenever $W_H(c)\ne0$
in an extension field,
the coefficients of $T^0,\ldots,T^{q-1}$ in $H(c+T)$ span its full
joint coefficient-matrix space.
\end{lemma}
\begin{proof}
Choose a basis $C_1,\ldots,C_q$ of the coefficient-matrix space and
write $H(T)=\sum_{\nu=1}^q f_\nu(T)C_\nu$. The scalar coordinate
polynomials $f_\nu$ are linearly independent; otherwise the original
coefficient matrices would span a proper subspace of this basis.

Let $D^{[j]}f(C)$ denote the coefficient of $T^j$ in $f(C+T)$.
The Hasse-Wronskian
\[
 W_H(C)=\det\bigl(D^{[j]}f_\nu(C)\bigr)_{
                  0\le j<q,\ 1\le\nu\le q}
\]
is nonzero in the stated characteristic regime. Indeed, a constant
basis change makes the $f_\nu$ have distinct degrees
$d_1<\cdots<d_q\le d$. Its leading determinant coefficient, up to
the nonzero leading coefficients of the $f_\nu$, is
\[
 \det\left(\binom{d_\nu}{j}\right)_{j,\nu}
   =\frac{\prod_{\mu<\nu}(d_\nu-d_\mu)}
          {\prod_{j=0}^{q-1}j!}\ne0.
\]
Here $q\le d+1$, so every denominator and degree difference is nonzero
in the stated characteristic. Invertibility at $c$ proves the spanning
assertion. Translation itself preserves the full coefficient span.
\end{proof}

\begin{lemma}[From support to total degree]\label{lem:degree}
Suppose $H_i(Y_i)$ are polynomial matrices on disjoint singleton
variables and every coefficient of $H_i$ lies in the span of its
coefficients of degrees at most $\eta$. If $\prod_iH_i$ is concentrated
in support at most $t$, then its coefficient span is generated in
total degree at most $\eta t$.
\end{lemma}
\begin{proof}
Expand a spanning coefficient of support at most $t$ as a product of
local coefficients. Keep every unselected factor's constant coefficient
unchanged. In the selected factors, expand each coefficient using
degrees at most $\eta$. Disjointness makes every resulting product an
actual coefficient of total degree at most $\eta t$.
Concentration in total degree adds under multiplication of products on
disjoint variables by the same coefficient-product argument.
Every nonzero scalar entry of a concentrated matrix has a nonzero
coefficient among the spanning low-degree coefficients.
\end{proof}

\begin{lemma}[Moving a low-degree witness]\label{lem:move}
Let $p(Y)\ne0$ over a coefficient field $E$ have a nonzero homogeneous
component of degree at most $h$, and let $L$ be invertible and linear.
Then $p(LX)$ has a nonzero monomial of support at most $h$, and
$p(L\SV_h)\ne0$.
\end{lemma}
\begin{proof}
An invertible homogeneous linear change is a graded polynomial-ring
automorphism. It preserves the nonzeroness and degree of each
homogeneous component. A monomial of degree at most $h$ uses at most
$h$ variables.

For the last assertion, choose such a nonzero monomial. Set one
interpolation seed to the abscissa of each variable in its support,
and set unused amplitudes to zero. This restricts the original
coordinates to that support, with independent amplitudes in the
selected coordinates. The chosen coefficient survives. The argument
works over every extension coefficient field.
\end{proof}

\section{Proof of Theorem~\ref{thm:orbit}}

Use independent coordinates $Y=LX$. We construct shifts in the
\emph{original} $X$ coordinates, starting with ranks $w$ and descending.
All earlier seeds are retained symbolically as coefficient-field
elements.

\paragraph{Full-rank base case.}
Take $G_w=\SV_1$. For every nonzero linear form $\ell_i$,
\[
 \ell_i(\SV_1)=z_1p_i(y_1),\qquad
 p_i(y)=\sum_{j=1}^n L_{ij}\lambda_j(y)\ne0.
\]
The last inequality follows by evaluating at the abscissas: a nonzero
row of $L$ cannot give the zero interpolation polynomial.
Consequently $c_i+z_1p_i(y_1)$ is transcendental over the original
coefficient field. Every nonzero univariate determinant remains
nonzero at that argument. A full-rank interval consists of full-rank
factors, so its determinant, a product of such nonzero factors, survives.

In the characteristic regime of Lemma~\ref{lem:jets}, this same shift
also makes every nonzero local compound's Hasse-Wronskian nonzero:
apply the lemma to $H_{i,r}(T)=\bigwedge^r A_i(T)$, whose degree is at most
$r\delta$ and whose coefficient-space dimension is at most
$\binom wr^2$. Evaluating its nonzero Wronskian at
$c_i+\ell_i(G_w)$ cannot kill it.
These conditions persist after adding fresh shifts, because zeroing
their amplitudes recovers the preceding nonzero Wronskians.

\paragraph{Inductive coordinate bookkeeping.}
Suppose $G_{r+1}$ preserves every interval rank greater than $r$.
Let $E$ be its rational seed field. In independent residual coordinates
$Y$, replace each local factor by
\[
 \widetilde A_i(Y_i)=
 A_i\bigl(Y_i+c_i+\ell_i(G_{r+1})\bigr).
\]
These factors still use disjoint singleton $Y$ variables and have
degree at most $\delta$. Their generic ranks agree with the original
interval ranks, and their values at $Y=0$ are exactly the interval
values at the old original-coordinate shift.

Every local compound $\bigwedge^r\widetilde A_i$ has its full coefficient
span generated by degrees at most $\eta_r$. In the unrestricted
characteristic case this is its degree bound $r\delta$.
In the good-characteristic case, Lemma~\ref{lem:jets} and the maintained
Wronskian give degrees at most $q_{i,r}-1$, where
\[
 q_{i,r}-1\le
 \min\left\{r\delta,\binom wr^2-1\right\}.
\]
Identically zero compounds are trivial, and constant factors need only
their degree-zero coefficients.

For a minor $p(Y)$ of such a residual product, the polynomial
$p(LX)$ is the corresponding minor after the residual original
coordinate shift $X\mapsto X+G_{r+1}$. The affine constants have
already been absorbed above. The remaining substitution $Y=LX$
is homogeneous and invertible.

\paragraph{Local witnesses.}
Fix an interval of generic rank $r$. Greedily write it as
\[
 B_1\cdots B_mT,\qquad B_j=C_jA_{e_j},
\]
where each $B_j$ is a shortest rank-$r$ prefix and each $C_j$, and the
tail $T$, has rank greater than $r$ or is empty. Every interval within
a higher-rank run has rank greater than $r$ and is already preserved
at residual zero.

The assumed cone theorem restricts each such run to at most $b_r$
singleton $Y$ variables while preserving its cone.
Lemmas~\ref{lem:cone} and~\ref{lem:degree} concentrate its $r$-th
compound in total degree at most $\eta_r b_r$.
A local final factor's compound is concentrated in total degree at most
$\eta_r$. Composition on disjoint $Y$ variables gives
\[
\begin{array}{c|ccc}
 \text{product}&B_j&B_jB_{j+1}&B_mT\\ \hline
 \text{compound total-degree bound}
   &\eta_r(b_r+1)&\eta_r(2b_r+2)&\eta_r(2b_r+1).
\end{array}
\]
Each product has generic rank $r$, so it has a nonzero $r$-minor.
Such a minor has a nonzero homogeneous component of degree at most
\[
 h_r=\eta_r(2b_r+2).
\]
Lemma~\ref{lem:move} shows that a fresh $\SV_{h_r}$ in the residual
original $X$ coordinates preserves it. This holds for all local
witnesses in their common symbolic seed field.

\paragraph{Completing an interval.}
At this specialization write $B_j=U_jV_j$ through rank $r$.
Adjacent-pair preservation makes every $V_jU_{j+1}$ invertible, and
the separate $B_mT$ witness makes $V_mT$ full row rank. Therefore
\[
 U_1(V_1U_2)\cdots(V_{m-1}U_m)(V_mT)
\]
has rank $r$. The argument includes $m=1$ and an empty tail.

Set $G_r=G_{r+1}+\SV_{h_r}$. Previously preserved minors remain
nonzero because setting all new amplitudes to zero recovers the old
map. Specialization cannot increase rank. This proves the descending
induction simultaneously for every interval. Additivity gives
$G_1=\SV_K$.

For a scalar output, insert the constant matrices $e_1u$ and $ve_1^T$.
The augmented full product is $fE_{11}$, whose rank is one when
$f\ne0$. This proves scalar hitting with the same width and seed count.

\paragraph{Size and finite grid.}
Since $r\le w$ and $\delta\ge1$,
\[
 K\le\delta w\left(1+\sum_{r<w}(2b_r+2)\right).
\]
The recurrence for $D_w$ gives the previously established estimate
$k_{\rm geo}(w)\le4^{w-1}(w!)^2$ for $w\ge2$.
In the good-characteristic regime, $\eta_r\le\binom wr^2-1<4^w$
gives $K\le4^w k_{\rm geo}(w)$.
The scalar polynomial has total degree at most $n\delta$, including
after the affine change. Each SV coordinate has total seed degree
at most $n$. Thus $f(\SV_K)$ has degree at most $n^2\delta$.
A product grid with $n^2\delta+1$ values for each of its $2K$ seeds
hits this nonzero polynomial. These values may overlap the abscissas.

\section{Scope and comparison}

The cases $n=0$ or $\delta=0$ are constant-polynomial models and need
one evaluation; the displayed nontrivial parameter bound assumes
$n,\delta\ge1$. The proof is characteristic independent. The use of
an invertible $L$ is essential to Lemma~\ref{lem:move}; no statement
about singular affine projections or orbit closures follows.
A fixed common diagonal shift cannot replace the first SV pair for
unknown $L$: a nonzero linear form can vanish on that diagonal.
The Wronskian improvement is not asserted in small characteristic:
for example, the independent polynomials $1,T^p$ have a singular
matrix of their first two Taylor coefficients in characteristic $p$.

\begin{remark}[Elementary fallback]
The qualitative fixed-width conclusion does not require the geometric
input if one instead assumes the elementary rank-profile concentration
theorem from the main manuscript. For each $r$, let $W_r=\binom wr$,
$\mathcal R_r=\{\binom sr:r<s\le w\}$, and let $M_r$ be its stated
concentration bound for higher-rank runs after taking the $r$-th
compound. Replace $b_r$ everywhere in the proof by $M_r$; the cone lemma
is then unnecessary. The same argument uses
\[
 K_{\rm elem}=1+\sum_{r<w}\eta_r(2M_r+2).
\]
In good characteristic this is at most $4^w k_{\rm elem}(w)$, where
$k_{\rm elem}(w)=1+\sum_{r<w}(2M_r+2)=2^{O(w^2)}$ is the elementary
main-manuscript budget. Thus polynomial query complexity in $n,\delta$
for fixed $w$ already follows from the elementary input; the cone
theorem improves the width dependence to $2^{O(w\log w)}$.
\end{remark}

For comparison, Saha and Thankey give quasipolynomial hitting sets for
orbits of multilinear constant-width ROABPs~\cite{ST21}.
Bhargava and Ghosh give the stronger unknown-order orbit bound
$(ndw)^{O(w^2\log n\min\{w^2,d\log w\})}$ in characteristic zero or
greater than $d$~\cite{BG21}. At fixed width and fixed preimage individual
degree, these verified bounds are quasipolynomial, whereas the
cone-theorem consequence above is polynomial. In the stated
good-characteristic regime it is also polynomial in the preimage degree
for each fixed width. This comparison does not
assert that no later result exists.

\begin{thebibliography}{2}
\bibitem{ST21}
C. Saha and B. Thankey.
\emph{Hitting Sets for Orbits of Circuit Classes and Polynomial Families}.
APPROX/RANDOM 2021, Article 50, especially Theorem 8.
\url{https://doi.org/10.4230/LIPIcs.APPROX/RANDOM.2021.50}.
\bibitem{BG21}
V. Bhargava and S. Ghosh.
\emph{Improved Hitting Set for Orbit of ROABPs}.
APPROX/RANDOM 2021, Article 30.
\url{https://doi.org/10.4230/LIPIcs.APPROX/RANDOM.2021.30}.
\end{thebibliography}
\end{document}
